\documentclass[11pt, a4paper]{article}

% ---------------------------------------------------------------------------
% Packages
% ---------------------------------------------------------------------------
\usepackage[utf8]{inputenc}
\usepackage[T1]{fontenc}
\usepackage[english]{babel}
\usepackage[margin=2.5cm]{geometry}
\usepackage{times}
\usepackage{graphicx}
\usepackage{float}
\usepackage{enumitem}
\usepackage{amsmath}
\usepackage{setspace}
\usepackage[colorlinks=true, linkcolor=blue, urlcolor=blue, citecolor=blue]{hyperref}

\onehalfspacing

% ---------------------------------------------------------------------------
% Title page information
% ---------------------------------------------------------------------------
\title{
    \textbf{Sentinel} \\
    \large An Automated Electronic Climbing Belay Device \\
    \vspace{0.3cm}
    \large Project Proposal
}
\author{
    Pol Mas Fillat \\
    \small Norwegian University of Science and Technology (NTNU), Trondheim \\
    \small Universitat Politècnica de Catalunya (UPC), Barcelona \\
    \vspace{0.2cm}
    \small Supervisor: Amund Skavhaug
}
\date{\today}

\begin{document}

\maketitle
\thispagestyle{empty}

\vspace{1cm}

\begin{abstract}
\noindent
Sentinel is a mechatronic system designed to function as an automated
electronic climbing belay device -- an ``electric Grigri.'' The device
actively manages climbing rope slack during ascent and performs a
fast-acting emergency brake in the event of a fall, removing the need for a
human belayer. This document presents the initial project proposal for the
Master's Thesis, outlining the motivation, objectives, scope, and
requirements of the system to be developed on an ESP32-based embedded
platform.
\end{abstract}

\vspace{0.5cm}

\section{Introduction}

Rock and sport climbing rely fundamentally on the belay system: the
mechanism and human partner responsible for managing rope slack and
arresting a climber's fall. Traditional belaying requires a trained,
attentive human belayer holding the rope at all times, which introduces
risks related to human error, fatigue, distraction, and inconsistent
reaction times. Semi-automated mechanical devices, such as the Petzl
GriGri, assist with braking but still require a human operator to feed
rope and manage slack.

\begin{figure}[H]
    \centering
    \includegraphics[width=0.6\textwidth]{figures/petzl_grigri_reference.png}
    \caption{Reference CAD model of the Petzl GriGri mechanical assembly,
    used as a starting point for the design of Sentinel's rope-gripping
    mechanism.}
    \label{fig:petzl_grigri}
\end{figure}

The mechanical core of the Petzl GriGri, shown in Figure
\ref{fig:petzl_grigri}, is taken as the initial design reference for
Sentinel. Its camming/pinching mechanism, which clamps the rope against a
grooved wheel to generate braking friction, is a proven, compact, and
mechanically reliable solution for assisted-braking belay devices. Rather
than designing a rope-gripping mechanism from scratch, this project will
use the GriGri's mechanical layout as a baseline and extend it with motor
actuation, enabling the mechanism to be driven autonomously for both
slack management (feeding and taking in rope) and emergency braking,
rather than requiring a human to manually operate the cam lever and
control the rope by hand.

Sentinel proposes a fully automated, sensor-driven belay device: an
embedded mechatronic system capable of continuously monitoring rope
position, tension, and dynamic motion in real time, and of autonomously
regulating slack while remaining ready to trigger an emergency brake
within milliseconds of detecting a fall event. By combining a rotary
encoder, a load cell, and an inertial measurement unit (IMU) with a
motor-actuated braking and rope-management mechanism, the system aims to
provide a belay experience that is at least as safe and responsive as a
competent human belayer, while removing the dependency on continuous human
attention.

This project is developed as a Master's Thesis in Industrial Engineering
at the Norwegian University of Science and Technology (NTNU), Trondheim,
under the supervision of Amund Skavhaug, in collaboration with the
Universitat Politècnica de Catalunya (UPC).

\section{Main Objectives}

\begin{enumerate}[leftmargin=1.2cm]
    \item Design and build a functional mechatronic prototype of an
        automated belay device capable of managing rope slack during a
        climber's ascent without human intervention.
    \item Implement a real-time embedded control system, running on an
        ESP32 microcontroller, that fuses data from a rotary encoder, a
        load cell, and an IMU to estimate the climber's state (climbing,
        resting, falling).
    \item Develop and validate a fast-acting emergency braking mechanism
        capable of arresting a fall within a safe reaction time and
        stopping distance, comparable to or better than manual belay
        devices.
    \item Create a remote control device / smartphone app / smartwatch app to allow the climber to control the belay device.
    \item Characterize and document the mechanical, electrical, and
        control-software design through rigorous testing, in a form
        suitable for academic evaluation as a Master's Thesis.
\end{enumerate}

\section{Secondary Objectives}

\begin{enumerate}[leftmargin=1.2cm]
    \item Explore adaptive slack management strategies (e.g. dynamic rope
        feed-out/take-in) to improve climber comfort and reduce unwanted
        rope drag or sudden tension spikes.
    \item Investigate sensor fusion techniques (e.g. Kalman filtering)
        to robustly distinguish fall events from ordinary climbing
        dynamics (clipping, resting on the rope, lead falls).
    \item Design the system with modularity and extensibility in mind, so
        that future work could explore additional features such as
        wireless connectivity, data logging for training analysis, or
        integration with gym safety infrastructure.
    \item Evaluate the safety, reliability, and failure modes of the
        device against relevant climbing equipment standards (e.g. UIAA,
        EN 12841) as a reference framework, without necessarily seeking
        formal certification within the scope of this thesis.
\end{enumerate}

\section{Justification}

Climbing gyms and outdoor climbing activities are experiencing sustained
growth in popularity, and belaying remains one of the most safety-critical
tasks in the sport. Human belayer error -- including inattention, incorrect
technique, and short-rope or hard-catch incidents -- is a documented
contributor to climbing injuries and fatalities. While auto-belay devices
exist commercially (e.g. for top-rope-only gym use, using retractable
cable/webbing systems), there is currently no widely available solution
that automates belaying for lead climbing using a standard climbing rope
while actively managing slack in real time.

An automated smart belay device addresses this gap by combining
continuous, fatigue-free sensor monitoring with a mechanically robust
braking system. From an engineering standpoint, the project is justified
by the interdisciplinary challenge it poses: it requires the integration
of mechanical design (rope-handling mechanism, braking system), embedded
real-time software, sensor signal processing, and control theory, making
it a comprehensive and appropriate capstone project for a Master's Thesis
in Industrial Engineering.

\section{Scope}

The scope of this thesis project includes:

\begin{itemize}[leftmargin=1.2cm]
    \item Design and fabrication of a mechanical prototype capable of
        gripping, feeding, and braking a standard single climbing rope.
    \item Selection and integration of a rotary encoder, load cell, and
        IMU for real-time rope and motion sensing.
    \item Development of embedded C++ firmware on an ESP32 platform
        (using PlatformIO) implementing sensor acquisition, state
        estimation, slack-management control, and emergency-brake
        triggering logic.
    \item Bench and controlled-environment testing of the prototype,
        including simulated fall scenarios, to evaluate reaction time,
        braking distance, and false-trigger rate.
    \item Documentation of the design process, testing methodology, and
        results in the Master's Thesis document.
\end{itemize}

The following are explicitly \textbf{out of scope} for this thesis, unless
revisited as future work:

\begin{itemize}[leftmargin=1.2cm]
    \item Formal certification of the device against climbing safety
        standards (e.g. UIAA, CE/EN).
    \item Outdoor field testing with real climbers at height.
    \item Mass-manufacturability and industrial design for a commercial
        product.
\end{itemize}

\section{Requirements}

\subsection{Hardware Requirements}

\begin{itemize}[leftmargin=1.2cm]
    \item \textbf{Microcontroller:} ESP32-based development board,
        providing sufficient processing power and I/O for real-time
        sensor fusion and motor control.
    \item \textbf{Rotary encoder:} to measure rope displacement and
        feed/retrieval speed.
    \item \textbf{Load cell (with amplifier, e.g. HX711 or similar):} to
        measure rope tension, used both for slack detection and as an
        input to fall detection.
    \item \textbf{Inertial Measurement Unit (IMU):} to detect dynamic
        motion characteristic of a fall event (e.g. sudden acceleration
        changes).
    \item \textbf{Actuator:} a servo or BLDC motor mechanism responsible
        for rope feed/take-in and for actuating the mechanical brake.
    \item \textbf{Braking mechanism:} a mechanically robust, fail-safe
        braking element capable of arresting fall loads.
    \item \textbf{Power supply:} suitable battery and power management
        circuitry for a portable prototype.
    \item \textbf{Structural frame/housing:} to mount all components
        securely around the rope path.
\end{itemize}

\subsection{Software Requirements}

\begin{itemize}[leftmargin=1.2cm]
    \item \textbf{Firmware development environment:} C++ using
        PlatformIO, targeting the ESP32 (\texttt{esp32dev}) platform.
    \item \textbf{Real-time sensor acquisition:} drivers/routines for
        reading the rotary encoder, load cell, and IMU at rates
        sufficient for fall detection latency requirements.
    \item \textbf{State estimation / sensor fusion:} algorithms (e.g.
        thresholding, filtering, or Kalman-based fusion) to classify the
        climber's current state and detect fall onset reliably.
    \item \textbf{Motor control logic:} closed-loop control of the
        rope-management actuator for slack take-up/feed-out.
    \item \textbf{Emergency brake control:} deterministic, low-latency
        logic to trigger the braking mechanism upon fall detection.
    \item \textbf{Documentation and version control:} structured project
        documentation (\texttt{docs/}) and development log
        (\texttt{docs/DEV\_LOG.md}) tracking design decisions.
\end{itemize}

\section{Expected Contributions}

This thesis is expected to contribute:

\begin{enumerate}[leftmargin=1.2cm]
    \item A working proof-of-concept prototype demonstrating that
        automated, sensor-driven belaying with real-time slack
        management and fast emergency braking is technically feasible
        using accessible embedded hardware.
    \item An open, documented reference architecture (hardware and
        firmware) for future research or development of automated belay
        devices, including sensor fusion strategies for fall detection.
    \item Empirical performance data (reaction time, braking distance,
        false-positive/false-negative rates) characterizing the
        prototype's safety-relevant behavior under controlled test
        conditions.
    \item A foundation for future work extending the system toward
        higher reliability, certification readiness, and real-world
        climbing gym deployment.
\end{enumerate}

\end{document}
